\documentclass[11pt]{article}
\usepackage[a4paper,margin=25mm]{geometry}
\usepackage{amsmath,amssymb}
\usepackage[T1]{fontenc}
\usepackage{lmodern}
\usepackage{hyperref}
\hypersetup{colorlinks=true,urlcolor=blue}
\setlength{\parindent}{0pt}
\setlength{\parskip}{7pt}
\title{An unbounded family of cubic graphs with one prime cycle length}
\author{SucRunBug}
\date{1 October 2026}
\begin{document}
\maketitle
\textbf{Conjecture 00000000057 -- Verdict: FALSE.}

\textbf{Filed statement.} Every $n$-vertex graph with minimum degree at least $3$ contains $\Omega(\log n/\log\log n)$ cycles whose lengths are pairwise distinct primes.

\section*{Proof}
For every integer $m\ge1$, let $G_m$ be the disjoint union of $m$ copies of $K_4$. Equivalently, its vertex set is $\{0,\ldots,m-1\}\times\{0,1,2,3\}$, and two vertices are adjacent exactly when their first coordinates agree and their second coordinates differ.

The graph has $n=4m$ vertices and every vertex has degree $3$. Every simple cycle remains in one connected component, since every edge preserves the first coordinate. A simple cycle in that component uses at most its four vertices and at least three vertices. Its length is therefore either $3$ or $4$. Of these two lengths only $3$ is prime. Triangles exist, so the set of prime cycle lengths is exactly $\{3\}$, independently of $m$.

In particular, at most one cycle can be selected if the selected cycles must have pairwise distinct prime lengths. The number of triangles is irrelevant: all of them have the same length.

The claimed $\Omega$ bound means that some constants $a>0$ and $N_0$ work uniformly for all sufficiently large $n$ and all graphs satisfying the degree condition. But $\log(4m)/\log\log(4m)\to\infty$. Choose $m$ so large that $4m\ge N_0$ and $a\log(4m)/\log\log(4m)>1$. The graph $G_m$ then contradicts the claimed lower bound. This is an infinite counterexample family, rather than a small-graph test of an asymptotic statement.

\section*{Lean formalization}
The Lean project formalizes the family using \texttt{Fin m $\times$ Fin 4}, proves that every degree is $3$, and proves that any set of prime cycle lengths has cardinality at most one. It also proves $\log(4m)/\log\log(4m)\to\infty$ and, for every proposed positive uniform constant and vertex threshold, constructs a graph beyond that threshold that defeats the proposed lower bound. Thus the asymptotic obstruction is included in the formal proof.

\textbf{Reproduction.} In the supplied \texttt{lean4/} project, run
\texttt{lake exe cache get}, then \texttt{lake build} and
\texttt{lake env lean Check.lean}. The pinned versions are Lean and Mathlib
\texttt{v4.33.0}. The supplied \texttt{reproduce.py} performs independent
finite sanity checks; these checks are not the proof of any infinite assertion.

\textbf{Source.} \url{https://github.com/The-Last-Math-Competition/The-Last-Math-Competition/blob/main/conjectures/00000000057.md}
\end{document}
